\begin{codebox}
\codeline{\textcolor{cmtgray}{\#\ ==============================}}
\codeline{\textcolor{cmtgray}{\#\ 3.\ 对照：同一个"设备"概念}}
\codeline{\textcolor{cmtgray}{\#\ ==============================}}
\codeline{\textcolor{cmtgray}{\#\ 哲学视角的追问：}}
\codeline{\textcolor{cmtgray}{\#\ \ \ "设备"作为人造物（Artifact）的本质是什么？}}
\codeline{\textcolor{cmtgray}{\#\ \ \ 一台车床更换全部零件后还是原来那台车床吗？（忒修斯之船）}}
\codeblank{}
\codeline{\textcolor{cmtgray}{\#\ 工程视角的约定：}}
\codeline{Equipment\ \(\sqsubseteq\)\ PhysicalObject\ \ \ \ \ \ \ \ \ \ \ \ \ \ \ \#\ 设备是物理对象}
\codeline{Equipment\ \(\sqsubseteq\)\ \(\exists\)hasSerialNumber.string\ \ \ \ \ \ \#\ 每台设备有序列号（同一性判据）}
\codeline{CNCLathe\ \(\sqsubseteq\)\ CNCMachine\ \(\sqsubseteq\)\ Equipment\ \ \ \ \ \ \ \ \#\ 分类层次}
\codeline{Equipment\ \(\sqcap\)\ Material\ \(\sqsubseteq\)\ \(\bot\)\ \ \ \ \ \ \ \ \ \ \ \ \ \ \ \ \ \#\ 设备与材料不相交}
\codeblank{}
\codeline{\textcolor{cmtgray}{\#\ 工程本体论不回答"本质"问题，}}
\codeline{\textcolor{cmtgray}{\#\ 而是给出可检验、可推理、可共享的形式化约定：}}
\codeline{\textcolor{cmtgray}{\#\ 序列号就是这个信息系统中设备的同一性判据——问题被"约定"解决}}
\codeblank{}
\end{codebox}
\color{black}
